\scanpage{12}
\mspage{5}
appelé ici une ``maquette''.

\uline{Proposition} \inserted{Considérons la maquette d'une courbe \unsure{standard|standart} $X$.}
a) Soient $\alpha,\beta\in S$, alors $\alpha,\beta$ appartiennent à la $=$ composante conn. du graphe $G$, ssi. $X_\alpha$ et $X_\beta$ appartiennent à la $=$ comp. conn. de $X$. Donc \deleted{\unsure{l'on}} on a une bij. can.
\[
(11)\qquad\pi_0(G_X)\simeq\pi_0(X),
\]
en particulier $X$ est connexe ssi $G_X$ est connexe.

b) Supposons $X$ connexe i.e. $\widehat X$ connexe, i.e. $H^0(\widehat X,\underline O_{\widehat X})\unsure{\xleftarrow{\sim}|\simeq}k$ i.e dim $H^0(\widehat X,\underline O_{\widehat X})=1$.

Posons \deleted{$g(X)=\dim H^1(X,\underline O_X),\quad\chi(X)=\chi$}
\[
(12)\quad\left\{\begin{aligned}
g(\widehat X)&=\dim H^1(\widehat X,\underline O_{\widehat X})\\
\chi_{\mathrm{coh}}(\widehat X)&=\chi(X\deleted{\unsure{g}},\underline O_X)=1-g(X)
\end{aligned}\right.,
\]
on a alors
\[
(13)\quad\chi_{\mathrm{coh}}(\widehat X)
=\deleted{\unsure{\sum_{\alpha\in S}\chi_{\mathrm{coh}}(\widehat{\widetilde X}_\alpha)}}
\inserted{\chi(\widetilde{\widehat X})-\underbrace{\operatorname{card}A}_{\mu}}
=\sum_{\alpha\in S}(1-g_\alpha)-\mu
\]
i.e.
\[
(14)\quad\left\{\begin{aligned}
g(\widehat X)&=\sum g_\alpha+(-\operatorname{card}S+\underbrace{\operatorname{card}A}_{\mu}+1)\\
&=\sum g_\alpha+h_1
\end{aligned}\right.,
\]
\mspage{11}
